\documentclass[11pt]{article}
\usepackage[T1]{fontenc}
\usepackage{lmodern}
\usepackage[margin=1in]{geometry}
\usepackage{amsmath,amssymb,amsthm,mathtools}
\usepackage{microtype}
\usepackage[hidelinks]{hyperref}

\newtheorem{theorem}{Theorem}
\newtheorem{remark}{Remark}
\newcommand{\R}{\mathbb{R}}
\newcommand{\C}{\mathbb{C}}
\newcommand{\Hs}[1]{H^{#1}}
\newcommand{\eps}{\varepsilon}
\newcommand{\Lpinf}{L^\infty}

\title{E69: High-Regularity $L^\infty$ Bound for the\\Corrected Regularized Curvature Operator}
\author{Richard J. Reyes}
\date{October 26, 2026}

\begin{document}
\maketitle

\begin{abstract}
We prove that the corrected complex curvature operator is essentially bounded
whenever its input lies in $H^s(\R^n)$ with $s>n/2+2$ and the regularization
parameter $\eps$ is positive. The regularized reciprocal is bounded globally,
including at zeros of the field. The estimate is explicit up to the Sobolev
embedding constant and does not require the spatial dimension to be at most three.
\end{abstract}

\section{Statement}

For $\eps>0$ and $\alpha\in\R$, define, for $z\in\C$,
\begin{equation}\label{eq:reciprocal}
 R_\eps(z)
 := \frac{\overline z}{|z|^2+\eps^2\exp(-2\alpha|z|^2)},
\end{equation}
and, for a complex-valued field $\psi$,
\begin{equation}\label{eq:curvature}
 \Theta_\eps[\psi]
 := - (\Delta\psi)\,R_\eps(\psi)
 = -\Delta\psi\,
 \frac{\overline\psi}{|\psi|^2+\eps^2\exp(-2\alpha|\psi|^2)}.
\end{equation}

\begin{theorem}[E69: bounded regularized curvature]\label{thm:e69}
Let $n\ge1$, $\eps>0$, $\alpha\in\R$, and $s>n/2+2$.
If $\psi\in H^s(\R^n;\C)$, then
\[
 \Theta_\eps[\psi]\in L^\infty(\R^n;\C).
\]
In fact, with $\alpha_+:=\max\{\alpha,0\}$ and
\begin{equation}\label{eq:M}
 M_{\eps,\alpha}
 :=\max\left\{1,\frac{e^{2\alpha_+}}{\eps^2}\right\},
\end{equation}
there is a finite constant $C_{n,s}>0$ independent of
$\psi,\eps,\alpha$ such that
\begin{equation}\label{eq:main}
 \boxed{\quad
 \|\Theta_\eps[\psi]\|_{L^\infty(\R^n)}
 \le C_{n,s}M_{\eps,\alpha}\|\psi\|_{H^s(\R^n)}.
 \quad}
\end{equation}
\end{theorem}

\section{Proof}
\begin{proof}
We first estimate the Laplacian. With the unitary Fourier transform,
\[
 \widehat{\Delta\psi}(\xi)=-|\xi|^2\widehat\psi(\xi).
\]
Consequently,
\begin{align}
 \|\Delta\psi\|_{H^{s-2}}^2
 &=\int_{\R^n}(1+|\xi|^2)^{s-2}|\xi|^4
             |\widehat\psi(\xi)|^2\,d\xi\\
 &\le\int_{\R^n}(1+|\xi|^2)^s
             |\widehat\psi(\xi)|^2\,d\xi
  =\|\psi\|_{H^s}^2,
\end{align}
where we used $|\xi|^4\le(1+|\xi|^2)^2$. Thus
\begin{equation}\label{eq:lap-hs}
 \|\Delta\psi\|_{H^{s-2}}\le\|\psi\|_{H^s}.
\end{equation}

Since $s-2>n/2$, the Sobolev embedding
$H^{s-2}(\R^n)\hookrightarrow L^\infty(\R^n)$ applies. Therefore
\begin{equation}\label{eq:lap-linf}
 \|\Delta\psi\|_{L^\infty}
 \le C_{n,s}\|\Delta\psi\|_{H^{s-2}}
 \le C_{n,s}\|\psi\|_{H^s}.
\end{equation}
For completeness, the embedding may be obtained by Cauchy--Schwarz:
for $r=s-2>n/2$,
\begin{align}
 \|\widehat{\Delta\psi}\|_{L^1}
 &\le\left(\int_{\R^n}(1+|\xi|^2)^{-r}\,d\xi\right)^{1/2}
       \|\Delta\psi\|_{H^r},
\end{align}
and the integral is finite since its radial tail behaves as
$\int_1^\infty\rho^{n-1-2r}\,d\rho$. Fourier inversion then yields
\eqref{eq:lap-linf}.

Next we bound the reciprocal \eqref{eq:reciprocal}. For any $z\in\C$, the
denominator
\[
 D_\eps(z):=|z|^2+\eps^2e^{-2\alpha|z|^2}
\]
is strictly positive, because $\eps>0$. Put $t=|z|\ge0$.
If $0\le t\le1$, then $e^{-2\alpha t^2}\ge e^{-2\alpha_+}$, so
\[
 |R_\eps(z)|
 =\frac{t}{t^2+\eps^2e^{-2\alpha t^2}}
 \le\frac{e^{2\alpha_+}}{\eps^2}.
\]
If $t\ge1$, the denominator is at least $t^2$, and hence
\[
 |R_\eps(z)|\le\frac{t}{t^2}=\frac1t\le1.
\]
Thus, uniformly in $z$,
\begin{equation}\label{eq:rec-bound}
 \sup_{z\in\C}|R_\eps(z)|\le M_{\eps,\alpha}<\infty.
\end{equation}
In particular, $R_\eps(0)=0$, so nodes of $\psi$ produce no
singularity.

Finally, by \eqref{eq:curvature}, \eqref{eq:lap-linf}, and
\eqref{eq:rec-bound},
\begin{align*}
 \|\Theta_\eps[\psi]\|_{L^\infty}
 &\le\|\Delta\psi\|_{L^\infty}
       \|R_\eps(\psi)\|_{L^\infty}\\
 &\le C_{n,s}\|\psi\|_{H^s} M_{\eps,\alpha}.
\end{align*}
This proves \eqref{eq:main} and the asserted $L^\infty$ membership.
\end{proof}

\section{Scope and limitations}

\begin{remark}[Spatial dimensions]
The sufficient requirement $s>n/2+2$ applies in every finite spatial
dimension. It requires $s>5/2,3,7/2,4,9/2$ for $n=1,2,3,4,5$,
respectively. By contrast, the embedding $H^2\hookrightarrow L^\infty$
for $n\le3$ bounds $\psi$, not necessarily $\Delta\psi$;
$H^2$ regularity alone does not yield Theorem~\ref{thm:e69}.
\end{remark}

\begin{remark}[Other domains]
The same reasoning works on a flat torus, with the periodic Sobolev
embedding, and on a bounded Lipschitz extension domain $\Omega$
when an $H^s(\Omega)\to L^\infty(\Omega)$ embedding for second-order
derivatives is available (for example, under the usual suitable extension
and regularity hypotheses). In that case the constant depends on $\Omega$.
\end{remark}

\begin{remark}[Fixed regularization]
The bound is for fixed $\eps>0$ and does not imply a uniform estimate
as $\eps\downarrow0$. Nor does it establish existence, uniqueness, or
stability of a PDE flow involving $\Theta_\eps$.
\end{remark}

\begin{remark}[Choice of regularizer]
The claim uses the modulus-squared regularizer \eqref{eq:reciprocal}.
It does not extend automatically to the older complex denominator
$\psi+\eps e^{-\alpha|\psi|^2}$, which can vanish.
\end{remark}

\end{document}